

\subsection{Cavity count at fixed removed volume}
\label{sec:bunnybudget}

Does cooling improve because placement is better, or simply because more
material is removed? We separate these effects with 72 Bunny designs:
two exact volume budgets, four counts $n\in\{2,4,6,12\}$, three independent
design seeds, and \method{}, voxel-adjoint TD, and random feasible placement.
The budgets equal three or six radius-$0.05$ balls: $0.765\%$ or $1.529\%$
of the original Bunny volume. Within each allocation, every cavity has
radius $r_n=0.05(b/n)^{1/3}$, where $b\in\{3,6\}$ is the budget in ball-volume
units. Radii follow this constraint; only centers are optimized. All methods
share 1,024 candidates within each seed/allocation, retain outer clearance
$0.012$ and pairwise surface gap $0.014$, and start from the unperforated domain.
TWoS retains the 12-walk scout and independent 128-walk shortlist rescoring;
this ranking heuristic is not a finite-insertion confidence certificate.
The compiled voxel selector uses a finer $96\times112\times84$ FD/CG grid
than the original Bunny design test.



% The following measured paragraph is populated from the retained summaries.
At the full volume budget, increasing the TWoS count from two to twelve raises mean cooling from $37.58\%$ to $54.33\%$, while the radius decreases from $0.0721$ to $0.0397$. At twelve cavities, the voxel and random means are $55.02\%$ and $38.78\%$, respectively. The most useful allocation comparison is across budgets: twelve small TWoS cavities at half volume yield $47.26\%$ cooling (design range $47.03$--$47.67\%$), versus $47.45\%$ ($46.78$--$48.22\%$) for six larger cavities at full volume. The former also has $20.6\%$ less cavity surface area, by exact sphere geometry. This is similar measured cooling with less removal, not a proved equivalence or an advantage unique to TWoS. The voxel selector shows a similar allocation trend. A fourfold thinner WoS termination shell, $1.25\times10^{-4}$, changes cooling by at most $0.204$ percentage points for the pre-specified first-seed full-budget endpoint designs from both TD selectors. This sensitivity check does not bound the remaining bias.
An additional first-seed check of the half-volume comparison, chosen after inspecting the allocation results, gives $46.88\%$ versus $47.21\%$ cooling with the thinner shell, respectively.


All objective comparisons use four independently scrambled 512-point
interior blocks and 128 WoS paths per active point, separate from selection.
Removed material contributes zero under fixed original-volume normalization;
reported percentages average paired block ratios to the initial state.
Design ranges and conditional block standard errors are retained separately;
three seeds do not establish equivalence or a globally optimal allocation.

\paragraph{Why a preferred count needs a cost.}
Additional nested ideal cooled removal lowers compliance exactly
(\cref{prop:coolingmonotonicity}), whereas different equal-volume allocations
are not nested. At fixed removed volume $V$, exposed area grows as $n^{1/3}$
and the sum of isolated capacities as $n^{2/3}$ (\cref{prop:budgetgeometry});
the latter is not the interacting union capacity. A volume-only penalty is
constant within this sweep. Instead, \cref{fig:bunnybudget}c illustrates
$J/J_0+\kappa n$ for an explicitly assumed count cost. This isolates the
modeling choice behind a ``sweet spot'' without attributing it to the
topological estimator or claiming an application-specific optimum.
